%% ma: L12-B08-0127
%% lop: 12
%% chuong: 2
%% bai: 8
%% dang: 12.8.24
%% loai: TN4
%% muc: VD
%% dap_an: "C"
%% nguon: "CĐ2.3.pdf giáo viên gửi 27/10/2026 (M07, Lê Minh Khôi – Nguyễn Anh Khôi, Chương II Chủ đề 3 Biểu thức tọa độ của các phép toán vectơ), Câu 42"
%% kien_thuc: "Bài 8 (tích vô hướng, điều kiện vuông góc), lớp 10 (phương trình bậc hai)"
%% trang_thai: chinh-thuc
%% ghi_chu: "Giáo viên đã duyệt 27/10/2026. Đã sửa lỗi in phương án C (thiếu căn của 2). Đề gốc in phương án C có nghiệm thứ hai (-căn 26 + 2)/6 (thiếu căn); đã sửa thành (-căn 26 + căn 2)/6"
\begin{ex}
Trong không gian $Oxyz$, cho hai vectơ $\vec a=(2;1;-2)$, $\vec b=\left(0;-\sqrt2;\sqrt2\right)$. Tìm tất cả giá trị của tham số $m$ để hai vectơ $\vec u=2\vec a+3m\vec b$ và $\vec v=m\vec a-\vec b$ vuông góc với nhau.
\choice{$m=\dfrac{\sqrt{26}+\sqrt2}6$}{$m=\dfrac{-\sqrt{26}+\sqrt2}6$}{\True $m=\dfrac{\sqrt{26}+\sqrt2}6$ hoặc $m=\dfrac{-\sqrt{26}+\sqrt2}6$}{$\dfrac{-\sqrt{26}+\sqrt2}6\le m\le\dfrac{\sqrt{26}+\sqrt2}6$}
\loigiai{$|\vec a|^2=9$, $|\vec b|^2=4$, $\vec a\cdot\vec b=-3\sqrt2$.

$\vec u\cdot\vec v=2m\cdot9-2\vec a\cdot\vec b+3m^2\vec a\cdot\vec b-3m\cdot4=-9\sqrt2m^2+6m+6\sqrt2=0$

$\Leftrightarrow3\sqrt2m^2-2m-2\sqrt2=0\Leftrightarrow m=\dfrac{\sqrt2\pm\sqrt{26}}6$.}
\end{ex}
